\begin{proposition}[充分统计量外置与可逆外置的结构性不相容]\label{prop:pom-sufficient-statistic-residual-noninvertibility}
\leanverified{paper\_pom\_sufficient\_statistic\_residual\_noninvertibility}
固定 \(m\ge 1\) 与 \(x\in X_m\)。
在纤维 \(\Fold_m^{-1}(x)\) 上考虑逆代换次数随机变量 \(R_x=r(\omega)\)（定义 \ref{def:pom-fiber-rewrite-rv}）。
记纤维生成函数
\[
Z_x(t)=\sum_{\omega\in\Fold_m^{-1}(x)} t^{r(\omega)},
\qquad
d_m(x)=Z_x(1),
\qquad
d_x:=\deg Z_x.
\]
\par
设 \(\mathcal R:\Omega_m\to\{0,1\}^{\ast}\) 为任意外置残差，且仅依赖于充分统计量 \(R_x\)：存在函数 \(g\) 使得对所有 \(\omega\in\Fold_m^{-1}(x)\) 有
\(
\mathcal R(\omega)=g(r(\omega))
\)。
则映射
\[
\Fold_m^{-1}(x)\longrightarrow \{0,1\}^{\ast},
\qquad
\omega\longmapsto \mathcal R(\omega)
\]
的像集大小满足
\[
\card{\mathcal R(\Fold_m^{-1}(x))}\ \le\ d_x+1.
\]
因此当 \(d_m(x)>d_x+1\) 时，映射
\(
\omega\mapsto (x,\mathcal R(\omega))
\)
不可能在纤维上为单射，因而任何只外置 \(r(\omega)\) 的方案在该纤维上不可能可逆。
\par
并且在分量分解 \(\Gamma_x\simeq\bigsqcup_i P_{\ell_i}\) 下（定理 \ref{thm:pom-fiber-indset-factorization}），有闭式
\[
d_m(x)=\prod_i F_{\ell_i+2},
\qquad
d_x=\sum_i \left\lceil\frac{\ell_i}{2}\right\rceil,
\]
从而除去有限退化情形外一般有 \(d_m(x)\gg d_x\)。
\end{proposition}

\begin{corollary}[充分统计量外置的全局可逆覆盖比例上界]\label{cor:pom-sufficient-statistic-residual-global-coverage}
沿用命题 \ref{prop:pom-sufficient-statistic-residual-noninvertibility} 的记号。
设外置残差 \(\mathcal R:\Omega_m\to\{0,1\}^{\ast}\) 在每条纤维 \(\Fold_m^{-1}(x)\) 上仅依赖于充分统计量 \(r(\omega)\)。
令
\[
\Psi:\Omega_m\to X_m\times\{0,1\}^{\ast},
\qquad
\Psi(\omega):=\bigl(\Fold_m(\omega),\mathcal R(\omega)\bigr).
\]
则其像集大小满足
\[
\frac{|\Psi(\Omega_m)|}{2^m}
\le
\frac{1}{2^m}\sum_{x\in X_m}(d_x+1)
\le
\frac{(m+1)|X_m|}{2^m}.
\]
在稳定类型口径下 \(|X_m|=F_{m+2}=\Theta(\varphi^m)\)，因此
\[
\frac{|\Psi(\Omega_m)|}{2^m}
=
O\!\left(m\left(\frac{\varphi}{2}\right)^m\right),
\qquad m\to\infty,
\]
从而该类外置机制的可逆覆盖比例以指数速度趋于 \(0\)。
\end{corollary}
\begin{proof}
对每个 \(x\in X_m\)，命题 \ref{prop:pom-sufficient-statistic-residual-noninvertibility} 给出
\(
|\mathcal R(\Fold_m^{-1}(x))|\le d_x+1
\)。
因此 \(|\Psi(\Omega_m)|\le \sum_x(d_x+1)\)，两边除以 \(2^m\) 得第一条不等式。
另一方面，由 \(\Gamma_x\simeq\bigsqcup_i P_{\ell_i}\) 的分量分解，命题 \ref{prop:pom-sufficient-statistic-residual-noninvertibility} 给出
\(
d_x=\sum_i\lceil \ell_i/2\rceil\le \sum_i \ell_i=m
\)，
故 \(d_x+1\le m+1\) 对所有 \(x\) 成立，从而得到第二条不等式。
最后一式由 \(|X_m|=F_{m+2}=\Theta(\varphi^m)\) 代入得到。证毕。
\end{proof}

\begin{corollary}[纤维多重度的 Fibonacci 素因子封闭性]\label{cor:pom-fiber-multiplicity-fibonacci-prime-closure}
沿用命题 \ref{prop:pom-sufficient-statistic-residual-noninvertibility} 的分量分解记号。
则对任意 \(x\in X_m\) 与任意素数 \(p\)，若 \(p\mid d_m(x)\)，则存在某个 \(r\le m+2\) 使得 \(p\mid F_r\)。
等价地，纤维多重度的全部素因子都来自有限初段 Fibonacci 数列 \(\{F_1,\dots,F_{m+2}\}\)。
\end{corollary}
\leanverified{paper\_pom\_fiber\_multiplicity\_fibonacci\_prime\_closure}
\begin{proof}
由 \(d_m(x)=\prod_i F_{\ell_i+2}\) 立得：若 \(p\mid d_m(x)\)，则 \(p\) 整除某一因子 \(F_{\ell_i+2}\)。
而每个分量长度满足 \(\ell_i\le m\)，故 \(\ell_i+2\le m+2\)。
取 \(r=\ell_i+2\) 即得。证毕。
\end{proof}

\begin{proof}
由于 \(r(\omega)\in\{0,1,\dots,d_x\}\)，故 \(\mathcal R(\omega)=g(r(\omega))\) 的取值数至多为 \(d_x+1\)，得到 \(\card{\mathcal R(\Fold_m^{-1}(x))}\le d_x+1\)。
当 \(d_m(x)>d_x+1\) 时由抽屉原理必存在 \(\omega\neq\omega'\) 满足 \(\mathcal R(\omega)=\mathcal R(\omega')\)，从而 \(\omega\mapsto(x,\mathcal R(\omega))\) 不可能单射。
\par
最后给出闭式。
由定理 \ref{thm:pom-fiber-indset-factorization} 有 \(Z_x(t)=\prod_i I_{\ell_i}(t)\)，且 \(d_m(x)=Z_x(1)=\prod_i I_{\ell_i}(1)=\prod_i F_{\ell_i+2}\)。
另一方面，\(\deg I_{\ell}=\lceil \ell/2\rceil\)（最大独立集大小），并且在不交并下次数可加，
故 \(d_x=\deg Z_x=\sum_i \deg I_{\ell_i}=\sum_i \lceil \ell_i/2\rceil\)（亦见定理 \ref{thm:pom-fiber-reconstruction-cartesian-product} 的立方维数闭式）。
由 \(F_{\ell+2}\) 的指数增长而 \(\lceil \ell/2\rceil\) 仅线性增长，乘积与求和的尺度分离推出一般情形下 \(d_m(x)\gg d_x\)。
证毕。
\end{proof}

\begin{corollary}[充分统计外置与可逆外置不相容的完全退化分类]\label{cor:pom-sufficient-statistic-injective-degeneracy-classification}
\leanverified{paper\_pom\_sufficient\_statistic\_injective\_degeneracy\_classification}
沿用命题 \ref{prop:pom-sufficient-statistic-residual-noninvertibility} 的记号，并设 \(\mathcal R(\omega)=g(r(\omega))\) 在纤维 \(\Fold_m^{-1}(x)\) 上仅依赖充分统计 \(r(\omega)\)。
则除下列两类退化情形外，恒有严格不等式 \(d_m(x)>d_x+1\)（从而 \(\omega\mapsto(x,\mathcal R(\omega))\) 不可能在该纤维上单射）：
\begin{enumerate}
  \item \textup{若 \(d_x=0\)，则 \(Z_x\) 为常数多项式且 \(d_m(x)=1=d_x+1\)；}
  \item \textup{若 \(d_x=1\) 且 \(\Gamma_x\simeq P_1\)，则 \(d_m(x)=F_3=2=d_x+1\)。}
\end{enumerate}
更强地，对所有满足 \(d_x\ge 2\) 的稳定类型，成立指数下界
\[
\boxed{\ d_m(x)\ \ge\ 2^{d_x}\ >\ d_x+1\ }.
\]
\end{corollary}
\begin{proof}
先证引理：对一切整数 \(\ell\ge 1\) 有
\[
F_{\ell+2}\ \ge\ 2^{\lceil \ell/2\rceil}.
\]
对 \(\ell=1,2\) 直接检验成立。
若该不等式对 \(\ell,\ell+1\) 成立，则由 Fibonacci 递推 \(F_{\ell+4}=F_{\ell+3}+F_{\ell+2}\) 得
\[
F_{\ell+4}\ \ge\ 2^{\lceil (\ell+1)/2\rceil}+2^{\lceil \ell/2\rceil}
\ \ge\ 2^{\lceil (\ell+2)/2\rceil},
\]
其中最后一步由对 \(\ell\) 奇偶分类可直接验证；故引理对一切 \(\ell\) 成立。
\par
回到分量分解 \(\Gamma_x\simeq\bigsqcup_i P_{\ell_i}\)。
由命题 \ref{prop:pom-sufficient-statistic-residual-noninvertibility} 的闭式
\(
d_m(x)=\prod_i F_{\ell_i+2}
\) 与
\(
d_x=\sum_i\lceil \ell_i/2\rceil
\)，结合上述引理得到
\[
d_m(x)
=\prod_i F_{\ell_i+2}
\ \ge\ \prod_i 2^{\lceil \ell_i/2\rceil}
=2^{\sum_i\lceil \ell_i/2\rceil}
=2^{d_x}.
\]
当 \(d_x\ge 2\) 时有 \(2^{d_x}\ge 4>d_x+1\)，得到严格不等式。
\par
当 \(d_x=0\) 时必无分量，故 \(d_m(x)=1\)。
当 \(d_x=1\) 时仅存在单一分量且 \(\lceil \ell/2\rceil=1\) 强制 \(\ell\in\{1,2\}\)；
\(\ell=1\) 时 \(d_m(x)=F_3=2\) 恰取等，而 \(\ell=2\) 时 \(d_m(x)=F_4=3>2\)。
证毕。
\end{proof}

\begin{corollary}[充分统计外置与可逆外置之间的线性熵鸿沟]\label{cor:pom-sufficient-vs-invertible-externalization-entropy-gap}
\leanverified{paper\_pom\_sufficient\_vs\_invertible\_externalization\_entropy\_gap}
沿用命题 \ref{prop:pom-sufficient-statistic-residual-noninvertibility} 的分量分解 \(\Gamma_x\simeq\bigsqcup_{i=1}^c P_{\ell_i}\)，并记总顶点数
\[
L:=\sum_{i=1}^{c}\ell_i.
\]
则有下界
\[
\boxed{\;
\log d_m(x)-\log(d_x+1)
\ \ge\
(\log\varphi)\,L\ -\ 2c\log\varphi\ -\ \log(L+c+2)\; }.
\]
特别地，当分量数 \(c\) 在某一类稳定类型上被一致上界时，上式给出一个主项为 \(L\log\varphi\) 的线性鸿沟，而修正项至多为 \(O(\log L)\)。
\end{corollary}
\begin{proof}
由标准 Fibonacci 下界 \(F_n\ge \varphi^{n-2}\)（\(n\ge 2\)）与 \(d_m(x)=\prod_i F_{\ell_i+2}\)，得
\[
\log d_m(x)=\sum_{i=1}^{c}\log F_{\ell_i+2}
\ \ge\ \sum_{i=1}^{c}\ell_i\log\varphi-2c\log\varphi
=L\log\varphi-2c\log\varphi.
\]
另一方面 \(d_x=\sum_i\lceil \ell_i/2\rceil\le (L+c)/2\)，故
\(
\log(d_x+1)\le \log((L+c)/2+1)\le \log(L+c+2)
\)。
两式相减即得结论。证毕。
\end{proof}

\endinput
